\begin{textsample}{POS Dim 1 – vem\_ed\_26 – Score 6.00 – vem\_ed\_26/t137.txt}  \label{ex:f1_pos_062}
Leitores Aos Osvaldo Succi Jr.
Coordenador PCIs O mês de outubro de 2024 \textbf{foi} repleto de eventos e atividades internacionais.
Entre os dias 21 e 24, ocorreu a IVEC 2024, principal conferência internacional sobre Intercâmbios Virtuais – neste ano, totalmente on-line.
A equipe dos Projetos Colaborativos Internacionais (PCIs / Cesu) participou ativamente, com cinco trabalhos: duas discussões facilitadas, uma mesa redonda e duas"Flash Sessions"assíncronas.
Na segunda quinzena de outubro, a equipe dos PCIs também esteve envolvida nas visitas internacionais do Kirkwood Community College (EUA) à Fatec Itaquera e da Unisimón (Colômbia) à Fatec São Paulo.
O reconhecimento pelo \textbf{número} de PCIs desenvolvidos (cerca de 120 ao ano), pela diversidade de temas abordados e pela participação em eventos internacionais \textbf{resultou} em um convite para minha \textbf{integração} à diretoria executiva da COIL Connect. \textbf{É} uma organização sem fins lucrativos voltada ao fomento da Internacionalização em Casa no Ensino Superior, por meio dos projetos COIL (Collaborative Online International Learning).
Na seção"Artigo de Opinião", a professora Maria Edna Gomes e suas orientandas na Fatec Barueri discutem a \textbf{importância} dos PCIs como política \textbf{linguística} para a formação dos \textbf{futuros} tecnólogos, com \textbf{foco} no PCI realizado em língua espanhola no curso de Comércio Exterior da unidade em parceria com a DUOC UC (Chile).
Quer contribuir você também?
Acesse https://publicacoescesu.cps.sp.gov.br/VEm/about/submissions.
Dúvidas?
Escreva para cesu.pci@cps.sp.gov.br Boa leitura!

% matched lemmas: foco, futuro, importância, integração, linguístico, número, resultar, ser
\end{textsample}
